%%%%%%%%%%%%%%%%%%%%%%%%%%%%%%%%%%%%%%%%%%%%%%%%%%%%%%%%%%%%%%%%%%%%%%%%%%%%%%%%%%%%%%%%%%%%%%%%%%%%%%%%%%%%%%%%%%%%%%%%%%%%%%%%%%%%%%%%%%%%%%%%%%%%%%%%%%%
% This is just an example/guide for you to refer to when submitting manuscripts to Frontiers, it is not mandatory to use Frontiers .cls files nor frontiers.tex  %
% This will only generate the Manuscript, the final article will be typeset by Frontiers after acceptance.   
%                                              %
%                                                                                                                                                         %
% When submitting your files, remember to upload this *tex file, the pdf generated with it, the *bib file (if bibliography is not within the *tex) and all the figures.
%%%%%%%%%%%%%%%%%%%%%%%%%%%%%%%%%%%%%%%%%%%%%%%%%%%%%%%%%%%%%%%%%%%%%%%%%%%%%%%%%%%%%%%%%%%%%%%%%%%%%%%%%%%%%%%%%%%%%%%%%%%%%%%%%%%%%%%%%%%%%%%%%%%%%%%%%%%

%%% Version 3.4 Generated 2018/06/15 %%%
%%% You will need to have the following packages installed: datetime, fmtcount, etoolbox, fcprefix, which are normally inlcuded in WinEdt. %%%
%%% In http://www.ctan.org/ you can find the packages and how to install them, if necessary. %%%
%%%  NB logo1.jpg is required in the path in order to correctly compile front page header %%%

\documentclass[utf8]{frontiersSCNS} % for Science, Engineering and Humanities and Social Sciences articles
%\documentclass[utf8]{frontiersHLTH} % for Health articles
%\documentclass[utf8]{frontiersFPHY} % for Physics and Applied Mathematics and Statistics articles

%\setcitestyle{square} % for Physics and Applied Mathematics and Statistics articles
\usepackage{url,hyperref,lineno,microtype,subcaption}
\usepackage[onehalfspacing]{setspace}

\linenumbers


% Leave a blank line between paragraphs instead of using \\


\def\keyFont{\fontsize{8}{11}\helveticabold }
\def\firstAuthorLast{Clayton} %use et al only if is more than 1 author
\def\Authors{Sophie Clayton$^{1,*}$}
% Affiliations should be keyed to the author's name with superscript numbers and be listed as follows: Laboratory, Institute, Department, Organization, City, State abbreviation (USA, Canada, Australia), and Country (without detailed address information such as city zip codes or street names).
% If one of the authors has a change of address, list the new address below the correspondence details using a superscript symbol and use the same symbol to indicate the author in the author list.
\def\Address{$^{1}$Department of Ocean and Earth Sciences, Old Dominion University, Norfolk , VA, USA}
% The Corresponding Author should be marked with an asterisk
% Provide the exact contact address (this time including street name and city zip code) and email of the corresponding author
\def\corrAuthor{Sophie Clayton}

\def\corrEmail{sclayton@odu.edu}




\begin{document}
\onecolumn
\firstpage{1}

\title{Spatial and temporal variability in nutricline depth and nitrate fluxes in the Kuroshio and Oyashio Extensions from BGC-Argo profiles
} 

\author[\firstAuthorLast]{\Authors} %This field will be automatically populated
\address{} %This field will be automatically populated
\correspondance{} %This field will be automatically populated

\maketitle


\begin{abstract}
The depth of the nutricline and the magnitude of vertical nutrient fluxes have a direct influence on primary productivity and the strength of the biological carbon pump. Western boundary currents have been shown to be hotspots for carbon cycling, largely due to their high levels of eddy kinetic energy superimposed over large latitudinal gradients in physical and biogeochemical properties. The presence of mesoscale eddies and submesoscale filaments drives variability in the mixed layer and the nutricline depth over a wide range of space and time scales. The Kuroshio and Oyashio have been shown to act as large-scale subsurface nutrient streams, supporting large lateral transports of nutrients within the upper thermocline, and these nutrient streams are characterized by persistent along-isopycnal positive nitrate anomalies. However the extent of the spatial and temporal variability of the nutricline depth in the north western Pacific is not well known, partly due to the relatively sparse number of nutrient profiles over the broader region. In this study, data from a large number of continuous nitrate profiles from BGC-Argo floats deployed in the western Pacific, spanning multiple years and seasons, are used to determine temporal and spatial variability in the nutricline depth and nitrate fluxes, and how they relate to physical dynamics and primary productivity estimated from remote sensing.


\section{}



\tiny
 \keyFont{ \section{Keywords:} keyword, keyword, keyword, keyword, keyword, keyword, keyword, keyword} %All article types: you may provide up to 8 keywords; at least 5 are mandatory.
\end{abstract}

\section{Introduction}

Nutricline estimates from continuous profiles rather than discrete data points

\section{Data and Methods}
\subsections{BGC-Argo float data}
List float numbers and duration of deployment
Mention additional QA/QC on float data \citep{maurer2021delayed}

Figure 1. Data density in time/space (map showing float trajectories and number of profiles in 5x5 dgree boxes, histogram showing prfolies by month of year)

\subsection{Nutricline estimation}
Which criterion to use? 

\subsection{Nitrate flux estimates}
Dissipation rate estimation using Thorpe scale method (refer to Whalen et al. 2012 and Han et a;., 2021)


section{Results}



\section{Discussion}


\section{Conclusions}


\section*{Conflict of Interest Statement}
%All financial, commercial or other relationships that might be perceived by the academic community as representing a potential conflict of interest must be disclosed. If no such relationship exists, authors will be asked to confirm the following statement: 

The author declares that the research was conducted in the absence of any commercial or financial relationships that could be construed as a potential conflict of interest.

\section*{Author Contributions}

The Author Contributions section is mandatory for all articles, including articles by sole authors. If an appropriate statement is not provided on submission, a standard one will be inserted during the production process. The Author Contributions statement must describe the contributions of individual authors referred to by their initials and, in doing so, all authors agree to be accountable for the content of the work. Please see  \href{http://home.frontiersin.org/about/author-guidelines#AuthorandContributors}{here} for full authorship criteria.

\section*{Funding}
LOL

\section*{Acknowledgments}


\bibliographystyle{frontiersinSCNS_ENG_HUMS} 

\bibliography{kuro}

%%% Make sure to upload the bib file along with the tex file and PDF
%%% Please see the test.bib file for some examples of references

\section*{Figure captions}

% only 2 figs/tables allowed

% \begin{figure}[h!]
% \begin{center}
% \includegraphics[width=10cm]{logo1}% This is a *.eps file
% \end{center}
% \caption{ Enter the caption for your figure here.  Repeat as  necessary for each of your figures}\label{fig:1}
% \end{figure}


% \begin{figure}[h!]
% \begin{center}
% \includegraphics[width=15cm]{logos}
% \end{center}
% \caption{This is a figure with sub figures, \textbf{(A)} is one logo, \textbf{(B)} is a different logo.}\label{fig:2}
% \end{figure}

%%% If you are submitting a figure with subfigures please combine these into one image file with part labels integrated.
%%% If you don't add the figures in the LaTeX files, please upload them when submitting the article.
%%% Frontiers will add the figures at the end of the provisional pdf automatically
%%% The use of LaTeX coding to draw Diagrams/Figures/Structures should be avoided. They should be external callouts including graphics.

\end{document}
